% sec3_3.tex — Section 3.3: The General Formulation

We now provide a general formulation.  It is a model described by the following set
of assumptions and is a generalization of the model of Chapter~2.

%% ─── Regressors and Instruments ────────────────────────────────────────────
\subsection*{Regressors and Instruments}

\begin{assumption}[linearity]
\label{ass:3.1}
\emph{The equation to be estimated is linear:}
\[
  y_i = \mathbf{z}_i' \boldsymbol{\delta} + \varepsilon_i
  \quad (i = 1, 2, \ldots, n),
\]
\emph{where $\mathbf{z}_i$ is an $L$-dimensional vector of regressors, $\boldsymbol{\delta}$ is an
$L$-dimensional coefficient vector, and $\varepsilon_i$ is an unobservable error term.}
\end{assumption}

\begin{assumption}[ergodic stationarity]
\label{ass:3.2}
\emph{Let $\mathbf{x}_i$ be a $K$-dimensional vector to be
referred to as the vector of instruments, and let $\mathbf{w}_i$ be the unique and nonconstant
elements of $(y_i, \mathbf{z}_i, \mathbf{x}_i)$.\footnote{See examples below for illustrations of
$\mathbf{w}_i$.}  $\{\mathbf{w}_i\}$ is jointly stationary and ergodic.}
\end{assumption}

\begin{assumption}[orthogonality conditions]
\label{ass:3.3}
\emph{All the $K$ variables in $\mathbf{x}_i$ are \textbf{predetermined} in the sense that they are all
orthogonal to the current error term:
$\E(x_{ik}\,\varepsilon_i) = 0$ for all $i$ and $k$ $(k = 1, 2, \ldots, K)$.\footnote{As was noted in
Section~2.3, our use of the term \emph{predetermined} may not be standard in some quarters
of the profession.  All we require for the term is that the current error term be orthogonal
to the current regressors.  We do not require that the current error term be orthogonal to
the past regressors.}  This can be written as}
\[
  \E[\mathbf{x}_i \cdot (y_i - \mathbf{z}_i' \boldsymbol{\delta})] = \mathbf{0}
  \quad \text{or} \quad
  \E(\mathbf{g}_i) = \mathbf{0},
\]
\emph{where $\mathbf{g}_i \equiv \mathbf{x}_i \cdot \varepsilon_i$.}
\end{assumption}

We remarked in Section~2.3 that, when the regressors include a constant, the orthogonality
conditions are equivalent to the condition that $\E(\varepsilon_i) = 0$ and that the
regressors be uncorrelated with the error term.  Here, the orthogonality conditions are
about the instrumental variables.  So if one of the instruments is a constant, Assumption~3.3
is equivalent to the condition that $\E(\varepsilon_i) = 0$ and that nonconstant instruments be
uncorrelated with the error term.

The examples examined in the previous two sections can be written in this
general format.  For example,

\begin{example}[Working's example with an observable supply shifter]
\label{ex:3.1}
Consider the market equilibrium model of \eqref{eq:3.1.2a} and \eqref{eq:3.1.2bp} on
page~189.  Suppose the equation to be estimated is the demand equation.  The example
can be cast in the general format by setting
\[
  y_i = q_i, \quad L = 2, \quad
  \boldsymbol{\delta} = \begin{bmatrix} \alpha_0 \\ \alpha_1 \end{bmatrix}, \quad
  \mathbf{z}_i = \begin{bmatrix} 1 \\ p_i \end{bmatrix}, \quad
  \varepsilon_i = u_i, \quad K = 2, \quad
  \mathbf{x}_i = \begin{bmatrix} 1 \\ x_i \end{bmatrix},
\]
and $\mathbf{w}_i = (q_i, p_i, x_i)'$.  Since the mean of the error term is zero, a constant is
orthogonal to the error term and so can be included in $\mathbf{x}_i$.  It follows that $x_i$,
which is assumed to be uncorrelated with the error term, satisfies the orthogonality
condition $\E(x_i \varepsilon_i) = 0$.  So it too can be included in $\mathbf{x}_i$.
\end{example}

The other examples of the previous sections can be similarly written as special
cases of the general model.

Having two separate symbols, $\mathbf{x}_i$ and $\mathbf{z}_i$, may give you the impression that
the regressors and the instruments do not share the same variables, but that is not
always the case.  Indeed in the above example, $\mathbf{x}_i$ and $\mathbf{z}_i$ share the same variable
(a constant).  The instruments that are also regressors are called \textbf{predetermined
regressors}, and the rest of the regressors, those that are not included in $\mathbf{x}_i$, are
called \textbf{endogenous regressors}.  A good example to make the point is

\begin{example}[wage equation]
\label{ex:3.2}
A simplified version of the wage equation to be estimated later in this chapter is
\[
  LW_i = \delta_1 + \delta_2 S_i + \delta_3 \mathit{EXPR}_i + \delta_4 \mathit{IQ}_i + \varepsilon_i,
\]
where $LW_i$ is the log wage of individual $i$, $S_i$ is completed years of schooling,
$\mathit{EXPR}_i$ is experience in years, $\mathit{IQ}_i$ is IQ, and $\varepsilon_i$ is unobservable individual
characteristics relevant for the wage rate.  We assume $\E(\varepsilon_i) = 0$ (if not,
include the mean of $\varepsilon_i$ in $\delta_1$).  In one of the specifications we estimate later, we
assume that $S_i$ is predetermined but that $\mathit{IQ}_i$, an error-ridden measure of the
individual's ability, is endogenous due to the errors-in-variables problem.  We
also assume $\mathit{EXPR}_i$, $\mathit{AGE}_i$ (age of the individual), and $\mathit{MED}_i$ (mother's
education in years) are predetermined.  $\mathit{AGE}_i$ is excluded from the wage equation,
reflecting the underlying assumption that, once experience is controlled
for, age has no effect on the wage rate.  In terms of the general model,
\[
  y_i = LW_i, \quad L = 4, \quad
  \mathbf{z}_i = \begin{bmatrix} 1 \\ S_i \\ \mathit{EXPR}_i \\ \mathit{IQ}_i \end{bmatrix}, \quad
  K = 5, \quad
  \mathbf{x}_i = \begin{bmatrix} 1 \\ S_i \\ \mathit{EXPR}_i \\ \mathit{AGE}_i \\ \mathit{MED}_i \end{bmatrix},
\]
and $\mathbf{w}_i = (LW_i, S_i, \mathit{EXPR}_i, \mathit{IQ}_i, \mathit{AGE}_i, \mathit{MED}_i)'$.
As in Example~3.1, we can include a constant in $\mathbf{x}_i$ because $\E(\varepsilon_i) = 0$.  In this
example, $\mathbf{x}_i$ and $\mathbf{z}_i$ share three variables $(1, S_i, \mathit{EXPR}_i)$.  If, for example,
$S_i$ were not included in $\mathbf{x}_i$, that would amount to treating $S_i$ as endogenous.
\end{example}

If, as in this example, some of the regressors $\mathbf{z}_i$ are predetermined, those predetermined
regressors should be included in the list of instruments $\mathbf{x}_i$.  The GMM
estimation of the parameter vector is about how to exploit the information afforded
by the orthogonality conditions.  Not including predetermined regressors as elements
of $\mathbf{x}_i$ amounts to throwing away the orthogonality conditions that could have
been exploited.

%% ─── Identification ────────────────────────────────────────────────────────
\subsection*{Identification}

As seen from the examples of the previous two sections, an instrument must be not
only predetermined (i.e., orthogonal to the error term) but also correlated with the
regressors.  Otherwise, the instrumental variables estimator cannot be defined (see,
e.g., \eqref{eq:3.1.11}).  The generalization to the case with more than one regressor and
more than one predetermined variable is

\begin{assumption}[rank condition for identification]
\label{ass:3.4}
\emph{The $K \times L$ matrix $\E(\mathbf{x}_i \mathbf{z}_i')$
is of full column rank (i.e., its rank equals $L$, the number of columns).  We
denote this matrix by $\boldsymbol{\Sigma}_{\mathbf{xz}}$.\footnote{So the cross moment
$\E(\mathbf{x}_i \mathbf{z}_i')$ is assumed to exist and is finite.  If a moment is indicated, as here,
then by implication the moment is assumed to exist and is finite.}}
\end{assumption}

To see that this is indeed a generalization, consider Example~3.1.  With
$\mathbf{z}_i = (1, p_i)'$, $\mathbf{x}_i = (1, x_i)'$, the $\boldsymbol{\Sigma}_{\mathbf{xz}}$
matrix is
\[
  \boldsymbol{\Sigma}_{\mathbf{xz}}
  = \begin{bmatrix} 1 & \E(p_i) \\ \E(x_i) & \E(x_i p_i) \end{bmatrix}.
\]
The determinant of this matrix is not zero (and hence is of full column rank) if and
only if $\Cov(x_i, p_i) = \E(x_i p_i) - \E(x_i)\,\E(p_i) \neq 0$.

Assumption~3.4 is called the \textbf{rank condition for identification} for the following
reason.  To emphasize the dependence of $\mathbf{g}_i$ ($\equiv \mathbf{x}_i \cdot \varepsilon_i$) on the data and the
parameter vector, rewrite $\mathbf{g}_i$ as
\begin{equation}
  \mathbf{g}_i = \mathbf{g}(\mathbf{w}_i;\, \boldsymbol{\delta})
  \equiv \mathbf{x}_i \cdot (y_i - \mathbf{z}_i' \boldsymbol{\delta}).
  \label{eq:3.3.1}
\end{equation}
So the orthogonality conditions can be rewritten as
\begin{equation}
  \E[\mathbf{g}(\mathbf{w}_i;\, \boldsymbol{\delta})]
  = \underset{(K \times 1)}{\mathbf{0}}.
  \label{eq:3.3.2}
\end{equation}
Now let $\tilde{\boldsymbol{\delta}}$ ($L \times 1$) be a hypothetical value of $\boldsymbol{\delta}$
and consider the system of $K$ simultaneous equations in $L$ unknowns (the $L$ elements
of $\tilde{\boldsymbol{\delta}}$):
\begin{equation}
  \E[\mathbf{g}(\mathbf{w}_i;\, \tilde{\boldsymbol{\delta}})]
  = \underset{(K \times 1)}{\mathbf{0}}.
  \label{eq:3.3.3}
\end{equation}
The orthogonality conditions \eqref{eq:3.3.2} mean that the true value of the coefficient
vector, $\boldsymbol{\delta}$, is \emph{a} solution to this system of $K$ simultaneous equations
\eqref{eq:3.3.3}.  So the assumptions we have made so far guarantee that there exists a
solution to the simultaneous equations system.  We say that the coefficient vector (or
the equation) is \textbf{identified} if $\tilde{\boldsymbol{\delta}} = \boldsymbol{\delta}$ is the
\emph{only} solution.

Because the estimation equation is linear in our model, the function
$\mathbf{g}(\mathbf{w}_i;\, \tilde{\boldsymbol{\delta}})$ is linear in $\tilde{\boldsymbol{\delta}}$,
as it can be written as
$\mathbf{x}_i \cdot y_i - \mathbf{x}_i \mathbf{z}_i' \tilde{\boldsymbol{\delta}}$.
So \eqref{eq:3.3.3} is a system of $K$ linear equations:
\begin{equation}
  \E(\mathbf{x}_i \cdot y_i) - \E(\mathbf{x}_i \mathbf{z}_i')
  \tilde{\boldsymbol{\delta}}
  = \mathbf{0}
  \quad \text{or} \quad
  \underset{(K \times L)}{\boldsymbol{\Sigma}_{\mathbf{xz}}} \;
  \underset{(L \times 1)}{\tilde{\boldsymbol{\delta}}}
  = \underset{(K \times 1)}{\boldsymbol{\sigma}_{\mathbf{xy}}},
  \label{eq:3.3.4}
\end{equation}
where
\[
  \boldsymbol{\sigma}_{\mathbf{xy}} \equiv \E(\mathbf{x}_i \cdot y_i), \quad
  \boldsymbol{\Sigma}_{\mathbf{xz}} \equiv \E(\mathbf{x}_i \mathbf{z}_i').
\]
A necessary and sufficient condition that $\tilde{\boldsymbol{\delta}} = \boldsymbol{\delta}$ is the only
solution can be derived from the following result from matrix algebra.\footnote{See, e.g.,
Searle (1982, pp.~233--234).}

Suppose there exists a solution to a system of linear simultaneous equations
in $\mathbf{x}$: $\mathbf{A}\mathbf{x} = \mathbf{b}$.  A necessary and sufficient condition that it is the only
solution is that $\mathbf{A}$ is of full column rank.

Therefore, $\tilde{\boldsymbol{\delta}} = \boldsymbol{\delta}$ is the only solution to \eqref{eq:3.3.4} if
and only if $\boldsymbol{\Sigma}_{\mathbf{xz}}$ is of full column rank, which is Assumption~3.4.

%% ─── Order Condition for Identification ────────────────────────────────────
\subsection*{Order Condition for Identification}

Since $\rank(\boldsymbol{\Sigma}_{\mathbf{xz}}) < L$ if $K < L$, a \emph{necessary}
condition for identification is that
\begin{equation}
  K \; (= \text{\#predetermined variables})
  \;\geq\; L \; (= \text{\#regressors}).
  \label{eq:3.3.5a}
\end{equation}
This is called the \textbf{order condition for identification}.  It can be stated in different
ways.  Since $K$ is also the number of orthogonality conditions and $L$ is the number
of parameters, the order condition can be stated equivalently as
\begin{equation}
  \text{\#orthogonality conditions} \;\geq\; \text{\#parameters}.
  \label{eq:3.3.5b}
\end{equation}
By subtracting the number of predetermined regressors from both sides of the
inequality, we obtain another equivalent statement:
\begin{equation}
  \text{\#predetermined variables excluded from the equation}
  \;\geq\; \text{\#endogenous regressors}.
  \label{eq:3.3.5c}
\end{equation}
Depending on whether the order condition is satisfied, we say that the equation is
\begin{itemize}
\item \textbf{overidentified} if the rank condition is satisfied and $K > L$,
\item \textbf{exactly identified} or \textbf{just identified} if the rank condition is satisfied
  and $K = L$,
\item \textbf{underidentified} (or not identified) if the order condition is not satisfied
  (i.e., if $K < L$).
\end{itemize}
Since the order condition is a necessary condition for identification, a failure of the
order condition means that the equation is not identified.

%% ─── The Assumption for Asymptotic Normality ──────────────────────────────
\subsection*{The Assumption for Asymptotic Normality}

As in Chapter~2, we need to strengthen the orthogonality conditions for the estimator
to be asymptotically normal.

\begin{assumption}[$\mathbf{g}_i$ is a martingale difference sequence with finite second
moments]
\label{ass:3.5}
\emph{Let $\mathbf{g}_i \equiv \mathbf{x}_i \cdot \varepsilon_i$.
$\{\mathbf{g}_i\}$ is a martingale difference sequence (so a fortiori
$\E(\mathbf{g}_i) = \mathbf{0}$).  The $K \times K$ matrix of cross moments,
$\E(\mathbf{g}_i \mathbf{g}_i')$, is nonsingular.
We use $\mathbf{S}$ for $\avar(\bar{\mathbf{g}})$ (i.e., the variance of the limiting distribution
of $\sqrt{n}\,\bar{\mathbf{g}}$, where
$\bar{\mathbf{g}} \equiv \frac{1}{n} \sum_{i=1}^{n} \mathbf{g}_i$).
By Assumption~3.2 and the ergodic stationary Martingale Differences CLT,
$\mathbf{S} = \E(\mathbf{g}_i \mathbf{g}_i')$.}
\end{assumption}

This is the same as Assumption~2.5, and the same comments apply:
\begin{itemize}
\item If the instruments include a constant, then this assumption implies that the error
  is a martingale difference sequence (and \emph{a fortiori} serially uncorrelated).

\item A sufficient and perhaps easier to understand condition for Assumption~3.5 is
  that
  \begin{equation}
    \E(\varepsilon_i \mid \varepsilon_{i-1}, \ldots, \varepsilon_1,
    \mathbf{x}_i, \mathbf{x}_{i-1}, \ldots, \mathbf{x}_1) = 0,
    \label{eq:3.3.6}
  \end{equation}
  which means that, besides being a martingale difference sequence, the error term
  is orthogonal not only to the current but also to the past instruments.

\item Since $\mathbf{g}_i \mathbf{g}_i' = \varepsilon_i^2\, \mathbf{x}_i \mathbf{x}_i'$,
  $\mathbf{S}$ is a matrix of fourth moments.  Consistent estimation of $\mathbf{S}$ will require a
  fourth-moment assumption to be specified in Assumption~3.6 below.

\item If $\{\mathbf{g}_i\}$ is serially correlated, then $\mathbf{S}$ (which is defined to be
  $\avar(\bar{\mathbf{g}})$) does not equal $\E(\mathbf{g}_i \mathbf{g}_i')$ and will take a more
  complicated form, as we will see in Chapter~6.
\end{itemize}

%% ─── Questions for Review ──────────────────────────────────────────────────
\bigskip
\begin{center}\rule{0.6\textwidth}{0.4pt}\end{center}
\noindent\textsc{Questions for Review}
\medskip

\begin{enumerate}
\item In the Working example of Example~3.1, is the demand equation identified?
  Overidentified?  Is the supply equation identified?

\item Suppose the rank condition is satisfied for the wage equation of Example~3.2.
  Is the equation overidentified?

\item In the production function example, no instrument other than a constant is specified.
  So $K = 1$ and $L = 2$.  The order condition informs us that the log output
  equation \eqref{eq:3.2.13} cannot be identified.  Write down the orthogonality condition
  and verify that there are infinitely many combinations of $(\phi_0, \phi_1)$ that satisfy
  the orthogonality condition.

\item Verify that the examples of Section~3.2 are special cases of the model of this
  section by specifying $(y_i, \mathbf{x}_i, \mathbf{z}_i)$ and writing down the rank condition for each
  example.
\end{enumerate}
